\section{Software, independent numerical checks and source attribution}
\subsection{A fixed scheme, not a collection of unrelated finite parts}
The software accepts the physical invariants in fixed external-label order.
For a two-point graph the sole invariant is $s=p_E^2$. For a three-point
graph they are $(s_1,s_2,s_3)=(p_{1,E}^2,p_{2,E}^2,p_{3,E}^2)$, with
$p_1+p_2+p_3=0$. The topology registry derives the inserted physical leg or
opposite physical leg from the generated graph itself.
The proper bubble retains the finite kinetic Taylor term $-s/6$; the
overall subtraction is then performed at fixed inserted-edge momentum.
The definitions in the preceding sections determine the finite scheme.
They must not be replaced by pole subtraction alone.

Shared input validation evaluates the binary64 Gram determinant using exact
rational arithmetic. A rounding tolerance is relative to the minimum of all
three Cauchy product scales, and therefore commutes with external-leg
permutations. An accepted near-boundary negative determinant is corrected by
moving the largest invariant inward until the effective binary64 Gram
determinant is nonnegative. Every kernel and analytic bound uses this same
effective physical point; the requested values, effective values and number
of correction steps are retained in the result. The larger of the first two
physical vectors is used as a reconstruction axis to avoid amplifying a
tiny-leg rounding error.

\subsection{Reference point and normalization}
At $M=1$, the two-point invariant is $s=1$ and the three-point invariants
are $s_1=s_2=s_3=1$. The following are unweighted scalar integrals in the
measure $d^6k/\pi^3\,d^6q/\pi^3$:
\begin{center}
\begin{tabular}{lr}
\toprule Topology & $J_R$\\\midrule
One-point graph & $0$ exactly\\
Overlapping two-point graph & $+0.00277755520593736$\\
Bubble-inserted two-point graph & $+0.000240153063424584$\\
Primitive $K_{2,3}$ three-point graph & $-0.0830063381643418$\\
Triangle-inserted three-point graph & $+0.0036834315027189$\\
Bubble-inserted three-point graph & $-0.00617095361238479$\\
\bottomrule
\end{tabular}
\end{center}
The three fixed-label placements in each insertion family coincide only at
this symmetric point. Asymmetric points are used to check the physical-leg
mapping. Graph weights $1/S$, vertex factors $(-g)^{E+2}$ and the ordinary
two-loop measure conversion $(4\pi)^{-6}$ remain separate.

\subsection{Independent numerical and upstream checks}
The primitive Gaussian effective scale and the nested-triangle effective
scale are checked against an independent exact spanning-tree/two-forest
construction of the graph polynomials $U,F$. The projective period
$\mathcal P=1/2$ of the primitive graph is integrated analytically and is
also reproduced numerically. Bubble insertion is checked against direct
six-dimensional angular averages, including a nonsymmetric physical point.
The overlap integral has a separate three-chart five-dimensional rational
Taylor-remainder oracle, distinct from its two-chart four-dimensional
logarithmic core.

The nested-triangle oracle first integrates the subtracted inner triangle
and then performs an outer momentum integral using a radial coordinate,
two angular coordinates and two inner simplex coordinates. Antithetic
$k\leftrightarrow-k$ samples reduce odd angular contributions. Independent
scrambled Sobol calculations use NumPy/SciPy only in the optional
verification environment. Per-scramble estimates, seeds, sample counts and
standard errors are retained. They are statistical diagnostics rather than
strict integration-error enclosures.

The official pySecDec~1.6.6 source was built in a separate environment and
used to generate and compile the original six-edge primitive graph in
$d=6-2\epsilon$. Twelve sectors remain after upstream symmetry finding.
QMC evaluations at the zero point, symmetric point and an asymmetric point
give a pole compatible with $1/(4\epsilon)$. Subtracting the zero-point
finite Laurent coefficient gives the finite BPH value for this primitive
graph, since it has no proper divergent subgraphs. In one recorded run,
the symmetric finite difference differs from the local core by
$5.37\times10^{-7}$ and the asymmetric difference by
$4.05\times10^{-7}$, meeting the recorded five-standard-deviation and
absolute-target criteria. These
are genuine upstream numerical runs; the bare finite coefficient by itself
is not called a BPH amplitude.

The task-local native build includes two syntax compatibility corrections
for the current Apple toolchain: FORM's historical \texttt{alignof} macro
does not override the C++11 keyword, and dependent SecDecUtil member
template calls supply an empty explicit template argument list. Neither
correction changes the arithmetic algorithm. The original source archive,
patch description, installed-file hashes and generated library hash are
recorded in the build and comparison JSON files.

The external feyncop source separately checks the generated graphs,
fixed-label weights and embedded divergent subgraphs. pySecDec's Python
graph interface checks all supported $U,F$ polynomials and dimensional
normalizations. None of these finite checks replaces the analytic arguments
in the preceding sections.

\subsection{Precision and reproducibility}
The core runs with Python's standard library. Optional array acceleration
and external integrators are kept separate. Quadrature order differences,
alternate coordinate maps and stochastic standard errors are explicitly
distinguished from rigorous error bounds. Analytic amplitude bounds use
outward numeric enclosures for the effective converted inputs. Inadequate
refinement or unsupported floating-point scales must be reported, rather
than inferred to be accurate from a finite return value.

The release records exact source hashes, output hashes, independently
executed tests and the freshness of each external check. A changed source
fingerprint makes previous external evidence stale until it is rerun. The
iteration ledger records substantive changes and feedback, rather than
counting repeated passing commands as new iterations. The authoritative
current counts and raw results accompany the software in \path{results/}.

\subsection{Sources and authorship}
The recursive framework follows A.~D.~Kennedy, \emph{A Simple Proof of the
BPH Theorem}, arXiv:hep-th/9612113, especially the treatment of recursive
subtraction, Taylor remainders and scale control. The explicit constants,
coordinate maps and implementation in this document were derived during
this collaboration; they are not attributed to that paper as published
results.

Modern context includes O.~Schnetz, \emph{$\phi^3$ theory at six loops},
arXiv:2505.15485, and G.~Heinrich et al., \emph{Numerical Scattering
Amplitudes with pySecDec}, arXiv:2305.19768. The former concerns minimal
subtraction renormalization functions, which are not numerically identical
to the massive finite Taylor scheme here. Exact primary-source links,
versions and the availability or execution status of published code are
retained in \path{docs/two-loop-current-sources.json} and
\path{docs/modern-references.json}.

Software, proofs and checks were produced in this collaboration. A student
report should distinguish independently understood and reproduced work,
software assistance and externally verified statements. This document
establishes the declared finite two-loop cases, not an arbitrary-loop
renormalization theorem, a massless limit or a Minkowski continuation.
